\begin{lemma}[A-priori $L^2(\gamma_n)$ bound on the obstacle multiplier]
Fix $n \in \mathbb{N}$.  Let $F \in L^2(\gamma_n)$ with $F \geq 0$
$\gamma_n$-almost everywhere, let $a \in \mathbb{R}$ with $a > 0$, and let
$u \in L^2(\gamma_n)$ satisfy
\begin{itemize}
  \item $u \in \mathcal{K}$, the obstacle cone of \noderef{ObstacleCone}, and
  \item $J(u) \leq J(w)$ for every $w \in \mathcal{K}$, with
    $J = $ the functional of \noderef{ObstacleFunctional} for this $F$ and $a$.
\end{itemize}
Then for every $\psi \in L^2(\gamma_n)$ of finite energy, i.e. with
$\mathcal{E}(\psi) \neq +\infty$ in the sense of \noderef{Energy},
\[
  \bigl| T(\psi) \bigr| \leq a\, \|\psi\|_{L^2(\gamma_n)} ,
  \qquad
  T(\psi) = \text{the multiplier } \noderef{ObstacleMultiplier} \text{ at } \psi .
\]

The bound is uniform in $\psi$: the same constant serves every $\psi$ in the form
domain $\{\mathcal{E} \neq +\infty\}$, with no dependence on $\psi$ beyond its
$L^2(\gamma_n)$ norm.  And the constant is $a$ itself, not a multiple of it --
there is no implicit $C$, and no factor depending on $n$, on $F$ or on $u$.  Every
measure here is the Gaussian measure $\gamma_n$ of \noderef{GaussMeasure}: the
norm is the $L^2(\gamma_n)$ norm, the almost-everywhere nonnegativity of $F$ is
$\gamma_n$-almost everywhere, and the integrals appearing in the proof are
$\gamma_n$-integrals.  All four hypotheses are used: $F \geq 0$ and $a > 0$ in
Step 5, and the cone membership and minimality of $u$ in Steps 3 and 4, through
\noderef{MultiplierPositive} and \noderef{LewyStampacchia} respectively.

This is the a priori estimate that lets the multiplier extend from the form domain
to a bounded linear functional on $L^2(\gamma_n)$, which is what makes a
representing function available at all.
\end{lemma}
\begin{proof}
\emph{Step 0: the multiplier as a difference of a form and an inner product.}
By \noderef{HermiteL2ZeroAEOne} the zeroth Hermite mode $h_0$ of
\noderef{HermiteL2} equals the constant $1$ $\gamma_n$-almost everywhere, so the
element $g = F - a\,h_0$ of $L^2(\gamma_n)$ satisfies $g = F - a$
$\gamma_n$-almost everywhere.  Hence, by \noderef{L2InnerGaussIntegral} rewriting
the load integral as an inner product,
\begin{equation*}
  T(\varphi) = \mathcal{B}(u,\varphi) - \langle g, \varphi\rangle
  \qquad\text{for every } \varphi \in L^2(\gamma_n), \tag{0}
\end{equation*}
with $\mathcal{B}$ the bilinear form of \noderef{EnergyBil}.

\emph{Step 1: $T$ is additive and homogeneous on the form domain.}
By \noderef{EnergyFormBilinear} the set $\{\mathcal{E} \neq +\infty\}$ is closed
under sums and scalar multiples, and on it $\mathcal{B}$ is symmetric and linear in
each argument.  Combining this with $(0)$ and the linearity of the inner product in
its second argument: if $\mathcal{E}(v), \mathcal{E}(w) \neq +\infty$ and
$r \in \mathbb{R}$ then
\begin{equation*}
  T(v+w) = T(v) + T(w), \qquad T(r\,v) = r\,T(v). \tag{1}
\end{equation*}

\emph{Step 2: the positive and negative parts of $\psi$ lie in the cone.}
Let $\psi^+ = \max(\psi,0)$ and $\psi^- = \max(-\psi,0)$.  Both are almost
everywhere strongly measurable and satisfy $|\psi^\pm| \leq |\psi|$ pointwise, so
both lie in $L^2(\gamma_n)$ by domination.  Both are nonnegative by construction.
For the energy, the maps $t \mapsto \max(t,0)$ and $t \mapsto \max(-t,0)$ are
$1$-Lipschitz on $\mathbb{R}$ and send $0$ to $0$, so \noderef{LatticeContraction}
applies and gives $\mathcal{E}(\psi^+) \neq +\infty$ and
$\mathcal{E}(\psi^-) \neq +\infty$ from $\mathcal{E}(\psi) \neq +\infty$.  By
\noderef{ObstacleCone} -- finite energy together with almost everywhere
nonnegativity -- we conclude $\psi^+, \psi^- \in \mathcal{K}$.

\emph{Step 3: $|T(\psi)| \leq T(|\psi|)$.}
Pointwise $\psi = \psi^+ - \psi^-$, hence $\psi = \psi^+ + (-1)\psi^-$ in
$L^2(\gamma_n)$, and $(1)$ gives
\[
  T(\psi) = T(\psi^+) - T(\psi^-).
\]
Since $u$ lies in the cone and minimises $J$ there, \noderef{MultiplierPositive}
applies to each of $\psi^+, \psi^- \in \mathcal{K}$ and gives
$T(\psi^+) \geq 0$ and $T(\psi^-) \geq 0$.  A difference of two nonnegative reals
has absolute value at most their sum, so
\[
  |T(\psi)| = |T(\psi^+) - T(\psi^-)| \leq T(\psi^+) + T(\psi^-)
    = T(\psi^+ + \psi^-),
\]
the last equality by $(1)$ again.  Pointwise $\psi^+ + \psi^- = |\psi|$, so this
is the claim.

\emph{Step 4: the Lewy--Stampacchia bound at $|\psi|$.}
The element $\psi^+ + \psi^-$ is nonnegative almost everywhere, being a sum of two
such, and has finite energy because the form domain is closed under sums
(\noderef{EnergyFormBilinear}); so $\psi^+ + \psi^- \in \mathcal{K}$ by
\noderef{ObstacleCone}.  The hypotheses on $F$, $a$ and $u$ are exactly those of
\noderef{LewyStampacchia}, which applied at this test element gives
\begin{equation*}
  T(\psi^+ + \psi^-)
    \leq \int_{\mathbb{R}^n} (a - F)^+ \,\bigl(\psi^+ + \psi^-\bigr) \, d\gamma_n
    = \int_{\mathbb{R}^n} (a - F)^+ \,|\psi| \, d\gamma_n , \tag{2}
\end{equation*}
writing $(a-F)^+ = \max(a-F, 0)$ and using $\psi^+ + \psi^- = |\psi|$ almost
everywhere.

\emph{Step 5: $0 \leq (a-F)^+ \leq a$ almost everywhere.}
Fix $x$ with $F(x) \geq 0$, which is almost every $x$.  If $a - F(x) \leq 0$ then
$(a - F(x))^+ = 0 \leq a$, using $a > 0$ for the second inequality.  If
$a - F(x) \geq 0$ then $(a-F(x))^+ = a - F(x) \leq a$, using $F(x) \geq 0$.  In
both cases $0 \leq (a - F(x))^+ \leq a$.  Multiplying by $|\psi(x)| \geq 0$,
\[
  (a-F(x))^+\,|\psi(x)| \leq a\,|\psi(x)|
  \qquad \text{for } \gamma_n\text{-almost every } x ,
\]
and both sides are $\gamma_n$-integrable: the left because $(a-F)^+$ lies in
$L^2(\gamma_n)$ by \noderef{PosPartConstSubMemLp} and the product of two
$L^2(\gamma_n)$ functions is integrable, and the right because $\psi \in
L^2(\gamma_n)$ is integrable against the probability measure $\gamma_n$
(\noderef{GaussIsProb}).  Monotonicity of the integral therefore gives
\begin{equation*}
  \int_{\mathbb{R}^n} (a-F)^+\,|\psi| \, d\gamma_n
    \leq \int_{\mathbb{R}^n} a\,|\psi| \, d\gamma_n
    = a \int_{\mathbb{R}^n} |\psi| \, d\gamma_n . \tag{3}
\end{equation*}

\emph{Step 6: from the $L^1$ integral to the $L^2$ norm.}
The $L^1(\gamma_n)$ norm of the image of $\psi$ under the inclusion of
\noderef{ToL1} is $\int |\psi| \, d\gamma_n$, since that inclusion does not change
the underlying function.  By \noderef{L2IntoL1},
$\|\psi\|_{L^1(\gamma_n)} \leq \|\psi\|_{L^2(\gamma_n)}$, so
\begin{equation*}
  \int_{\mathbb{R}^n} |\psi| \, d\gamma_n \leq \|\psi\|_{L^2(\gamma_n)} ,
  \qquad\text{hence}\qquad
  a \int_{\mathbb{R}^n} |\psi| \, d\gamma_n \leq a\,\|\psi\|_{L^2(\gamma_n)} ,
  \tag{4}
\end{equation*}
the last step multiplying by $a \geq 0$.

\emph{Conclusion.}
Chaining Step 3, $(2)$, $(3)$ and $(4)$,
\[
  |T(\psi)|
    \leq T(\psi^+ + \psi^-)
    \leq \int_{\mathbb{R}^n} (a-F)^+ |\psi| \, d\gamma_n
    \leq a \int_{\mathbb{R}^n} |\psi| \, d\gamma_n
    \leq a\,\|\psi\|_{L^2(\gamma_n)} .
\]
Since $\psi$ was an arbitrary element of the form domain, the bound holds uniformly
there with the stated constant $a$.
\end{proof}
